\chapter[Considerações finais]{CONSIDERAÇÕES FINAIS}

\section{Relação entre as disciplinas e a prática}

Considero que todas as disciplinas do curso contribuíram para minha atuação na \ac{AGES}, ainda que essa relação nem sempre fosse imediata. Mesmo os conteúdos mais teóricos, matemáticos e ligados aos fundamentos da computação ajudaram a construir meu raciocínio lógico. Ao longo dos projetos, passei a perceber melhor o valor dessa base para decompor problemas, compreender relações entre informações e analisar alternativas antes de implementar uma solução. Por isso, não associo a contribuição de uma disciplina apenas ao uso direto de uma tecnologia estudada em aula.

Entre as contribuições que reconheço com mais clareza estão as disciplinas ligadas à gestão e ao desenvolvimento ágil. Na matriz do curso, essa formação aparece em componentes como Processos de Software e Gerenciamento de Projetos de Software. O planejamento das \textit{sprints}, a organização das tarefas, as revisões de escopo e o acompanhamento dos \textit{squads} foram oportunidades de aproximar esses conhecimentos da prática. No ReadRace, por exemplo, a redução do tempo disponível para a \textit{Sprint} 2 exigiu rever os compromissos assumidos, enquanto a divisão em \textit{squads} tornou o acompanhamento dos colegas mais direto. Essas situações me ajudaram a compreender a importância de adaptar o planejamento às condições reais da equipe.

Também destaco Projeto e Arquitetura de Software, pela contribuição para compreender a organização das aplicações e as relações entre suas partes. Essa perspectiva se relaciona ao trabalho que realizei com \textit{backend}, modelagem de dados e revisão de integrações. Em diferentes projetos, precisei considerar como uma alteração afetaria outras funcionalidades, além de entender o trecho de código em que estava trabalhando. A experiência com banco de dados, desde a \ac{AGES} II, também aproximou minha prática dos conhecimentos de Banco de Dados I e II, especialmente na organização das entidades e de seus relacionamentos.

As quatro passagens pela \ac{AGES} permitiram reunir conhecimentos que, nas disciplinas, eram apresentados com focos distintos. Na execução de um projeto, requisitos, arquitetura, implementação, qualidade e gestão se relacionam continuamente. Essa vivência ampliou minha percepção sobre a formação em Engenharia de Software e sobre a importância de manter uma base ampla para lidar com problemas que não se limitam a uma única área do curso.

\section{Aprendizados ao longo das quatro AGES}

Quando iniciei minha trajetória na \ac{AGES}, ainda não tinha experiência com o desenvolvimento de software em projetos reais. Ao longo das quatro passagens, pude construir uma compreensão mais ampla desse trabalho, vivenciando diferentes processos e os desafios técnicos envolvidos em sua execução. O contato com levantamento de requisitos, planejamento, modelagem, implementação, testes, correções e entrega me permitiu acompanhar etapas do ciclo de vida de um software e entender como elas dependem umas das outras.

Um dos aprendizados mais importantes foi reconhecer o valor de planejar antes de executar. Compreender a necessidade do cliente, discutir as alternativas e organizar as tarefas passou a fazer mais sentido à medida que eu percebia o impacto dessas decisões no desenvolvimento. As experiências com dúvidas de requisitos, dependências entre funcionalidades e ajustes de escopo mostraram que a preparação faz parte da construção da solução e influencia o resultado do trabalho técnico.

Também aprendi, pela convivência com as equipes, quanto a comunicação interfere no sucesso de um projeto. Explicar uma tarefa com clareza, ouvir as dificuldades de um colega e alinhar expectativas são ações que ajudam cada integrante a contribuir melhor. Conforme assumi mais responsabilidades, percebi que acompanhar o trabalho exigia entender o que as pessoas precisavam para avançar e encontrar uma forma adequada de orientá-las. A experiência de desenvolver junto de colegas com diferentes conhecimentos tornou esse aprendizado especialmente significativo.

Essa trajetória ampliou minha visão da Engenharia de Software. Passei a compreender o desenvolvimento como um trabalho que reúne conhecimento técnico, planejamento e colaboração durante a evolução do produto. Hoje, reconheço a importância de dedicar atenção tanto à solução implementada quanto às condições em que a equipe a constrói, pois a qualidade da comunicação e da organização influencia diretamente o que conseguimos entregar juntos.

\section{Direção profissional}

Desde que ingressei em Engenharia de Software, meu interesse estava voltado a uma carreira mais técnica. A arquitetura de software era o caminho que eu desejava seguir, por isso as oportunidades de aprofundar meus conhecimentos de \textit{backend}, modelagem e organização das aplicações tinham especial importância para mim.

As experiências nas \ac{AGES} III e IV ampliaram essa perspectiva. O contato mais próximo com a organização das equipes e o acompanhamento das entregas me permitiu conhecer responsabilidades que passaram a fazer parte das minhas possibilidades de carreira. Ao mesmo tempo, a presença crescente da inteligência artificial no desenvolvimento de software me levou a refletir sobre o direcionamento que gostaria de dar à minha atuação. Considerando essas experiências e minha leitura desse contexto, vejo a gestão como a melhor alternativa profissional para mim neste momento.

Essa percepção ajuda a explicar a atenção que tenho dedicado ao desenvolvimento das minhas \textit{soft skills}. Conduzir as apresentações das \textit{Sprints} 2, 3 e 4 e acompanhar os colegas foram oportunidades de praticar habilidades que considero importantes para essa direção. Quero continuar aprimorando minha comunicação e minha capacidade de orientar o trabalho em equipe, aproveitando a base técnica construída ao longo do curso para compreender os desafios das pessoas e dos projetos que vier a acompanhar.

\section{Avaliação da experiência na AGES}

Ao considerar minhas quatro passagens pela \ac{AGES}, reconheço sua importância principalmente para quem chega à graduação com pouca ou nenhuma experiência no desenvolvimento de software em um projeto real. Esse era o meu caso. Trabalhar com demandas de clientes, prazos, decisões técnicas e entregas coletivas me permitiu compreender responsabilidades que eu ainda não conhecia apenas pelo estudo das disciplinas. Essa experiência teve grande impacto na conquista do meu primeiro estágio e também contribuiu para minha contratação posterior, pois eu já conseguia relacionar o que aprendia na universidade a situações concretas de desenvolvimento.

Outro aspecto que considero positivo é a composição das equipes com estudantes de diferentes etapas da formação. A convivência entre pessoas com conhecimentos e níveis de experiência distintos favorece a troca de aprendizado. Nas primeiras passagens, pude aprender com colegas mais experientes; nas seguintes, assumi mais responsabilidades de orientação e acompanhamento. Essa mudança de posição dentro da equipe ampliou minha compreensão sobre colaboração e sobre a importância de tornar o conhecimento acessível aos demais integrantes.

Como possibilidade de melhoria, considero que o tamanho das equipes e a complexidade dos projetos poderiam ser reavaliados diante das mudanças na forma de desenvolver software, especialmente com o uso de ferramentas de inteligência artificial. Na minha percepção, alguns projetos têm equipes maiores do que o necessário para o escopo proposto. Acredito que uma equipe de quatro a seis pessoas bem engajadas poderia atender a um projeto semelhante ao ReadRace, com uma distribuição mais direta das responsabilidades e mais espaço de participação para cada integrante.

Uma alternativa seria trabalhar com equipes menores; outra seria ampliar a complexidade dos projetos escolhidos quando houver mais pessoas envolvidas. O objetivo dessa sugestão é aproximar o dimensionamento do trabalho das oportunidades de aprendizado, para que todos tenham desafios compatíveis com sua etapa de formação e consigam acompanhar o impacto de suas contribuições. Essa avaliação não diminui o valor que a \ac{AGES} teve na minha trajetória: parte justamente da importância que atribuo à experiência e do interesse em que ela continue oferecendo uma participação significativa a cada estudante.
